SQL (Structured Query Language) is a declarative programming language that allows us to store, access, and manipulate data stored in databases. Each database contains tables, which are rectangular collections of data with rows and columns. This section gives a brief overview of the small subset of SQL used by CS 61A; the full language has many more features.

\subsection{Creating Tables}
\subsubsection{\lstinline{CREATE TABLE}}
Say we want to create a table \lstinline{lift} which has integer columns \lstinline{chair}, \lstinline{single}, and \lstinline{couple}. 
Each row should describe one chair and how many single people and couples are in that chair at a specific instance in time. 
To create a named table to represent this scenario, use \lstinline{CREATE TABLE} with column names and types:
\begin{lstlisting}
CREATE TABLE lift (
    chair INTEGER,
    single INTEGER,
    couple INTEGER
);
\end{lstlisting}

The table can then be populated with \lstinline{INSERT INTO} statements. For example:
\begin{lstlisting}
INSERT INTO lift VALUES (1, 4, 2);
\end{lstlisting}

We can also create a new table from a \lstinline{SELECT} statement using \lstinline{CREATE TABLE ... AS}:
\begin{lstlisting}
CREATE TABLE crowded_lifts AS
    SELECT * FROM lift WHERE single >= 5;
\end{lstlisting}

\subsection{Querying Tables}
The \lstinline{FROM} clause specifies the source table for a query. We can choose columns with \lstinline{SELECT}:
\begin{lstlisting}
SELECT name, birth_month FROM team;
\end{lstlisting}
An asterisk (\lstinline{*}) selects all columns. Arithmetic expressions can compute values from columns in each row. For example, if each couple has two people, this query computes the total number of people in each chair:
\begin{lstlisting}
SELECT chair, single + 2 * couple AS people FROM lift;
\end{lstlisting}
The \lstinline{AS} keyword names the resulting column.

\subsubsection{\lstinline{WHERE}}
\lstinline{WHERE} filters rows according to a boolean condition:
\begin{lstlisting}
SELECT name FROM team WHERE pet = "dog";
\end{lstlisting}
In SQL, \lstinline{=} checks equality; it does not assign a value.

\subsubsection{\lstinline{ORDER BY}}
\lstinline{ORDER BY} sorts result rows by a column. Add \lstinline{ASC} for ascending order or \lstinline{DESC} for descending order; ascending order is the default.
\begin{lstlisting}
SELECT name FROM team WHERE pet = "dog" ORDER BY name DESC;
\end{lstlisting}

\subsubsection{\lstinline{LIMIT}}
\lstinline{LIMIT} restricts the result to a maximum number of rows. For example:
\begin{lstlisting}
SELECT chair FROM lift ORDER BY chair LIMIT 3;
\end{lstlisting}

